\subsection{Kreuzprodukt in Koordinatendarstellung}
\label{subsec:vector_algebra_cross_product_coordinates}

Wir stellen $\vec a,\vec b\in\R^3$ mithilfe der kanonischen Einheitsvektoren
dar:
\begin{align*}
    \vec a&=a_1\vec e_1+a_2\vec e_2+a_3\vec e_3 \\
    \vec b&=b_1\vec e_1+b_2\vec e_2+b_3\vec e_3.
\end{align*}
%
Mit den Kreuzprodukten der kanonischen Einheitsvektoren folgt durch
Ausmultiplizieren
\begin{ClozeArea}[34]{16}
  \[
    \begin{aligned}
      \vec a\times\vec b
      &=(a_1\vec e_1+a_2\vec e_2+a_3\vec e_3)
        \times(b_1\vec e_1+b_2\vec e_2+b_3\vec e_3)\\[1.5em]
      &=a_1b_1\overbrace{(\vec e_1\times\vec e_1)}^{=\,\vec 0}
        +a_1b_2\overbrace{(\vec e_1\times\vec e_2)}^{=\,\vec e_3}
        +a_1b_3\overbrace{(\vec e_1\times\vec e_3)}^{=\,-\vec e_2}
        \\[0.9em]
      &\quad
        +a_2b_1\overbrace{(\vec e_2\times\vec e_1)}^{=\,-\vec e_3}
        +a_2b_2\overbrace{(\vec e_2\times\vec e_2)}^{=\,\vec 0}
        +a_2b_3\overbrace{(\vec e_2\times\vec e_3)}^{=\,\vec e_1}
        \\[0.9em]
      &\quad
        +a_3b_1\overbrace{(\vec e_3\times\vec e_1)}^{=\,\vec e_2}
        +a_3b_2\overbrace{(\vec e_3\times\vec e_2)}^{=\,-\vec e_1}
        +a_3b_3\overbrace{(\vec e_3\times\vec e_3)}^{=\,\vec 0}
        \\[1.3em]
      &=(a_2b_3-a_3b_2)\vec e_1
        +(a_3b_1-a_1b_3)\vec e_2
        +(a_1b_2-a_2b_1)\vec e_3\\[1em]
      &=\begin{pmatrix}
          a_2b_3-a_3b_2\\
          a_3b_1-a_1b_3\\
          a_1b_2-a_2b_1
        \end{pmatrix}.
    \end{aligned}
  \]
\end{ClozeArea}

\begin{TheoremBox}{Koordinatenformel}{vector_algebra_cross_product_coordinate_formula}
  Für das Kreuzprodukt der Vektoren $\vec a,\vec b\in\R^3$ in Koordinatenform gilt
  \[
    \vec a\times\vec b
    =\begin{ClozeContinuation}[22]{3}
      \begin{pmatrix} a_1 \\ a_2 \\ a_3 \end{pmatrix}
      \times \begin{pmatrix} b_1 \\ b_2 \\ b_3 \end{pmatrix}
      =
      \begin{pmatrix} a_2b_3-a_3b_2 \\ a_3b_1-a_1b_3 \\ a_1b_2-a_2b_1 \end{pmatrix}
    \end{ClozeContinuation}.
  \]
\end{TheoremBox}

\begin{ExampleBox}{Kreuzprodukt aus Koordinaten}{vector_algebra_cross_product_coordinates_example}
  Für die folgenden Vektoren
  \[
    \vec a=\begin{pmatrix}2\\1\\-2\end{pmatrix},
    \qquad
    \vec b=\begin{pmatrix}-1\\2\\-1\end{pmatrix}
  \]
  ergibt sich das Kreuzprodukt zu
  \[
    \vec c=\vec a\times\vec b
    =\begin{ClozeContinuation}[26]{3}
      \begin{pmatrix}2\\1\\-2\end{pmatrix}
      \times\begin{pmatrix}-1\\2\\-1\end{pmatrix}
      =
      \begin{pmatrix}
        1\cdot(-1)-(-2)\cdot2\\
        (-2)\cdot(-1)-2\cdot(-1)\\
        2\cdot2-1\cdot(-1)
      \end{pmatrix}
      =\begin{pmatrix}3\\4\\5\end{pmatrix}
    \end{ClozeContinuation}.
  \]
  Eine sinnvolle Kontrolle des berechneten Kreuzprodukts ist die Prüfung
  seiner Orthogonalität zu beiden Ausgangsvektoren:
  \[
    \begin{aligned}
      \vec c\cdot\vec a
      &=\begin{ClozeContinuation}[24]{3}
        \begin{pmatrix}3\\4\\5\end{pmatrix}
        \cdot\begin{pmatrix}2\\1\\-2\end{pmatrix}
        =3\cdot2+4\cdot1+5\cdot(-2)=0
      \end{ClozeContinuation}
      \IfClozeTF{}{,\quad\raisebox{-0.15ex}{%
        \CachedTikzFigure{content/resources/chapters/vector_algebra/tikz/verified_checkmark.tex}}}\\[1em]
      \vec c\cdot\vec b
      &=\begin{ClozeContinuation}[24]{3}
        \begin{pmatrix}3\\4\\5\end{pmatrix}
        \cdot\begin{pmatrix}-1\\2\\-1\end{pmatrix}
        =3\cdot(-1)+4\cdot2+5\cdot(-1)=0
      \end{ClozeContinuation}
      \IfClozeTF{}{.\quad\raisebox{-0.15ex}{%
        \CachedTikzFigure{content/resources/chapters/vector_algebra/tikz/verified_checkmark.tex}}}
    \end{aligned}
  \]
  Der aufgespannte Flächeninhalt beträgt
  \[
    A=|\vec a\times\vec b|
    =\Cloze{7cm}{\sqrt{3^2+4^2+5^2}=\sqrt{50}}.
  \]
\end{ExampleBox}

Aus der Koordinatenformel in
Satz~\ref{thm:vector_algebra_cross_product_coordinate_formula} gewinnen wir
die Determinantenschreibweise schrittweise. Zunächst zerlegen wir den
Koordinatenvektor in die kanonischen Einheitsvektoren:
\begin{ClozeArea}[34]{12}
  \[
    \begin{aligned}
      \vec a\times\vec b
      &=\begin{pmatrix}
          a_2b_3-a_3b_2\\
          a_3b_1-a_1b_3\\
          a_1b_2-a_2b_1
        \end{pmatrix}
        =(a_2b_3-a_3b_2)\vec e_1
        +(a_3b_1-a_1b_3)\vec e_2
        +(a_1b_2-a_2b_1)\vec e_3
        \\[1.4em]
      &=\det\!\begin{pmatrix}
          a_2&b_2\\
          a_3&b_3
        \end{pmatrix}\vec e_1
        -\det\!\begin{pmatrix}
          a_1&b_1\\
          a_3&b_3
        \end{pmatrix}\vec e_2
        +\det\!\begin{pmatrix}
          a_1&b_1\\
          a_2&b_2
        \end{pmatrix}\vec e_3
        \\[1.4em]
      &=\det\!\begin{pmatrix}
          \MathHighlight{\vec e_1}&a_1&b_1\\
          \MathHighlight{\vec e_2}&a_2&b_2\\
          \MathHighlight{\vec e_3}&a_3&b_3
        \end{pmatrix}.
    \end{aligned}
  \]
\end{ClozeArea}
Die letzte Gleichheit liest die Entwicklung nach der gelb markierten ersten
Spalte rückwärts. Diese Schreibweise dient nur als formale Merkhilfe;
definiert wurde das Kreuzprodukt geometrisch.

\begin{KeyPointBox}{Determinantenschreibweise}{vector_algebra_cross_product_determinant_notation}
  Das Kreuzprodukt lässt sich formal schreiben als
  \[
    \vec a\times\vec b
    =\begin{ClozeContinuation}[29]{4}
      \det\!\begin{pmatrix}
        \MathHighlight{\vec e_1}&a_1&b_1 \\
        \MathHighlight{\vec e_2}&a_2&b_2 \\
        \MathHighlight{\vec e_3}&a_3&b_3
      \end{pmatrix}
      =\vec e_1\det\!\begin{pmatrix}a_2&b_2\\a_3&b_3\end{pmatrix}
      -\vec e_2\det\!\begin{pmatrix}a_1&b_1\\a_3&b_3\end{pmatrix}
      +\vec e_3\det\!\begin{pmatrix}a_1&b_1\\a_2&b_2\end{pmatrix}
    \end{ClozeContinuation}.
  \]
\end{KeyPointBox}

Für die Berechnung ist daher keine neue Regel erforderlich: Wir entwickeln die
Determinante mit der bereits bekannten Laplace-Entwicklung und berechnen die
entstehenden $2\times2$-Determinanten.

\begin{ExampleBox}{Kreuzprodukt in Determinantenschreibweise}{vector_algebra_cross_product_determinant_example}
  Wir berechnen das Kreuzprodukt aus
  Beispiel~\ref{ex:vector_algebra_cross_product_coordinates_example} erneut
  für
  \[
    \vec a=\begin{pmatrix}2\\1\\-2\end{pmatrix},
    \qquad
    \vec b=\begin{pmatrix}-1\\2\\-1\end{pmatrix},
  \]
  nun jedoch mithilfe der Determinantenschreibweise:
  \begin{ClozeArea}[32]{11}
    \[
      \begin{aligned}
        \vec a\times\vec b
        &=\det\!\begin{pmatrix}
            \MathHighlight{\vec e_1}&2&-1\\
            \MathHighlight{\vec e_2}&1& 2\\
            \MathHighlight{\vec e_3}&-2&-1
          \end{pmatrix}\\[1em]
        &=\vec e_1\det\!\begin{pmatrix}1&2\\-2&-1\end{pmatrix}
          -\vec e_2\det\!\begin{pmatrix}2&-1\\-2&-1\end{pmatrix}
          +\vec e_3\det\!\begin{pmatrix}2&-1\\1&2\end{pmatrix}\\[1em]
        &=\vec e_1\bigl(1\cdot(-1)-2\cdot(-2)\bigr)
          -\vec e_2\bigl(2\cdot(-1)-(-1)\cdot(-2)\bigr)
          +\vec e_3\bigl(2\cdot2-(-1)\cdot1\bigr)  \\[1em]
        &=3\vec e_1+4\vec e_2+5\vec e_3
          =\begin{pmatrix}3\\4\\5\end{pmatrix}.
      \end{aligned}
    \]
  \end{ClozeArea}
\end{ExampleBox}
